\chapter{Working like an expert}
\label{ch:expert}

You now know every part of \nodwyn. This last chapter is about habits: how
people who use it every day set it up, work with it, and keep it honest.

\section{Set up a project once}

\begin{enumerate}
  \item \textbf{Standing approvals for the boring, safe commands}: tests,
    build, lint, \code{git status}, \code{git diff}. In
    \file{.nodwyn/allow.toml}, so they travel with the repository.
  \item \textbf{A few sharp tools} in \file{.nodwyn/init.scm}: how this
    project runs its tests, how it builds, one or two readers that know the
    project's layout (like \code{read-adr}).
  \item \textbf{Teach the evidence loop your test command}, if it is not a
    common one (\code{test-command-words}).
  \item \textbf{Project facts in the system prompt}, by wrapping
    \code{system-prompt}: what not to touch, how things are done here.
  \item \textbf{Deny checks for the hard rules}: generated folders, package
    installs, whatever must never happen here.
  \item \textbf{Commit it}, so everyone can review the same setup. Each person
    grants trust separately with \code{nodwyn trust --repo PATH}, then
    \code{--accept REVISION}; changes to either authority file need another review.
\end{enumerate}

\section{Work in small, reviewable turns}

\begin{itemize}
  \item Name the file, the pattern to follow and the test to run. A local
    model does well with a paragraph and badly with a page.
  \item Read the diffs as they go by: every write is shown with its diff,
    whether or not it asked. An \code{fs-write} of a whole existing file
    deserves suspicion; an \code{fs-edit} is usually small.
  \item For a larger change: commit first, and review with \code{git diff}
    after. A write you do not want is one \code{git checkout} away for a
    committed file, and one \code{nodwyn undo} or Restore away for any
    other.
  \item When a turn goes sideways, stop it (\key{Ctrl-C}, or Stop in the
    console) and say what you want differently, or send the correction as
    a steer (\key{Ctrl+S} in the terminal client, Steer in the console),
    which stops the turn and sends it next. The log keeps what
    happened.
\end{itemize}

\section{Keep it honest}

\begin{itemize}
  \item \textbf{Read the reasons.} Every confirm tells you what a yes would
    do. The day you stop reading them is the day one matters.
  \item \textbf{Read the log when surprised.} \code{jq} and a
    \code{guard.decision} answer most ``why'' questions in a minute.
  \item \textbf{Prefer deny checks to switches.} A strategy can only skip; a
    check is enforced.
  \item \textbf{Run \code{/optimize} now and then}, with a focus, and review
    its diff like a colleague's.
  \item \textbf{Pair only devices you would hand a terminal}, and revoke
    the ones you no longer use: \code{nodwyn devices revoke NAME}.
\end{itemize}

\section{Where to go from here}

\begin{itemize}
  \item \file{docs/scheme.md} is the language's contract: every form and
    primitive, exactly.
  \item \file{internal/loop/*.scm} is the default loop. Read
    \file{strategy.scm}, then \file{agents.scm}, then the habits.
  \item \file{internal/core/userland\_test.go} has larger strategies written
    the way a user would write them: a plan cycle with a keyed judge, and
    routing between models.
  \item \file{docs/adr/} explains why things are the way they are. When you
    want to change something fundamental, start there.
\end{itemize}

And when you write something good, a habit, a tool, a check, share it: it
is a Scheme file, and the next person can \code{load} it.
